\section{Unit Testing of the Original Code}
\label{sec:unit-testing}

\subsection{Requirements Compliance}
\todo{Go through each requirement in the exam specification and state whether
the original code fulfils it. Refer to line numbers in
\code{ExplodingKittens.java}.}

\begin{longtable}{@{}p{1.2cm}p{6.5cm}p{1.8cm}p{5cm}@{}}
\toprule
\textbf{Req.} & \textbf{Description} & \textbf{Fulfilled?} & \textbf{Comment / line ref.} \\
\midrule
\endhead
\req{1} & \todo{e.g. 2--5 players, mix of bots and online players} & \todo{Yes/No/Partly} & \todo{} \\
\req{2} & \todo{e.g. deck composition depends on number of players} & \todo{} & \todo{} \\
\req{3} & \todo{} & \todo{} & \todo{} \\
\req{4} & \todo{} & \todo{} & \todo{} \\
\req{5} & \todo{} & \todo{} & \todo{} \\
\bottomrule
\caption{Requirements compliance of the original implementation.}
\label{tab:req-compliance}
\end{longtable}

\subsection{Testability of the Original Code}
\todo{Explain why the code is hard to unit test: everything in one class,
static mutable state (\code{deck}, \code{discard},
\code{numberOfTurnsToTake}), I/O mixed with game logic, \code{System.exit},
hard-coded \code{Random}, threads/timeouts, no dependency injection.}

\subsection{Identified Bugs}
\todo{List defects found while analysing/testing, e.g.:}
\begin{itemize}
  \item \todo{Two-of-a-kind uses \code{rnd.nextInt(size-1)} --- last card can never be stolen / crashes on a one-card hand.}
  \item \todo{Bots never act (no bot logic implemented).}
  \item \todo{Unvalidated input (target IDs, defuse position, card names) throws exceptions.}
  \item \todo{Exploded players can be targeted.}
  \item \todo{\ldots}
\end{itemize}
